\uuid{TaET}
\exo7id{2958}
\titre{exo7 2958}
\auteur{quercia}
\organisation{exo7}
\datecreate{2010-03-08}
\isIndication{false}
\isCorrection{true}
\chapitre{Nombres complexes}
\sousChapitre{Trigonométrie}
\cpge{mpsi-mp2i.nombres-complexes.c, pcsi.nombres-complexes.c, ptsi.nombres-complexes.c, mpsi-mp2i.calcul-algebrique-trigonometrie.d, pcsi.calcul-algebrique-trigonometrie.d, ptsi.calcul-algebrique-trigonometrie.d, mpsi-mp2i.calcul-algebrique-trigonometrie.a, pcsi.calcul-algebrique-trigonometrie.a, ptsi.calcul-algebrique-trigonometrie.a, bcpst1.trigonometrie, bcpst1.nombres-complexes.b}
\module{Algèbre}
\niveau{L1}
\difficulte{}

\contenu{
\question{
Soit $n \in \N$, et $x \in \R$.
Exprimer $\tan(nx)$ en fonction de $\tan x$.
}
\reponse{
$\tan(nx) = \frac {C_n^1\tan x - C_n^3\tan^3 x + \cdots}
                            {C_n^0       - C_n^2\tan^2 x + \cdots}$.
}
}
